\chapter{System Implementation and Technical Realization}
\label{ch:implementation}

This chapter describes the technical implementation of each subsystem in the proposed framework. The system is implemented as a set of modular, containerized services. Each service is developed and tested independently before being connected into a complete reporting pipeline. The implementation follows the architectural design described in Chapter~\ref{ch:architecture}.

\section{Private Permissioned Blockchain Infrastructure}
\label{sec:blockchain-impl}

The blockchain layer is the foundation of the proposed framework. It provides an immutable, tamper-evident ledger for recording civic reports, community votes, authority actions, and governance decisions.

\subsection{Network Setup and Configuration}

The blockchain network is implemented using Go Ethereum (Geth), version 1.13.15. The Clique \ac{PoA} consensus algorithm is used. Clique is a lightweight consensus mechanism designed for private networks. It assigns block-sealing authority to a set of pre-approved validator accounts. No mining or staking is required.

The network is initialized using a custom \texttt{genesis.json} file. This file defines the chain identifier, the initial authorized sealers, and the block period. The block time is set to five seconds. This provides deterministic and fast block finality, which is essential for a responsive civic reporting system.

\subsection{Node Topology}

The network consists of four nodes deployed on virtual private server (VPS) infrastructure:

\begin{enumerate}
    \item \textbf{Validator Node 1 — Municipal Authority}: Represents local government infrastructure departments. This node participates in block sealing.
    \item \textbf{Validator Node 2 — Regional Oversight}: Represents a regional or provincial supervisory body. This node participates in block sealing.
    \item \textbf{Validator Node 3 — Civic NGO}: Represents an independent civil society organization. This node participates in block sealing and prevents unilateral government censorship.
    \item \textbf{RPC Node (Node 4)}: A non-validator node that exposes a secure JSON-RPC endpoint. All off-chain services (the backend relayer, the DApp) connect to the blockchain through this node only. This separates the consensus layer from the application-facing layer.
\end{enumerate}

Each node is assigned a unique cryptographic identity. The addresses of the three validator nodes are recorded in the genesis block as authorized sealers. The peer connections between nodes are established using static node configuration (\texttt{static-nodes.json}) to ensure a stable and private network topology. The entire network operates within a private subnet and is not connected to any public Ethereum network.

Figure~\ref{fig:validator-nodes} (see Chapter~\ref{ch:architecture}) shows the 3-validator network topology.

\section{Smart Contract Architecture}
\label{sec:smart-contracts-impl}

Three Solidity smart contracts are deployed on the private PoA blockchain. They are compiled using Solidity version 0.8.20 and the Hardhat development framework. OpenZeppelin libraries are used for access control and security utilities.

\subsection{Reporting.sol}

\texttt{Reporting.sol} is the primary smart contract. It manages the entire lifecycle of civic issue reports.

\subsubsection*{Report Data Structure}

Each civic report is stored in a \texttt{Report} struct on-chain. The struct contains:

\begin{itemize}
    \item A unique integer \texttt{reportId}.
    \item The IPFS URI (\texttt{ipfsUri}) pointing to the off-chain report content.
    \item The payload hash (\texttt{payloadHash}) computed from the report description and ticket.
    \item The citizen's derived pseudonym (\texttt{citizenPseudonym}).
    \item Voting counters for validation, verification, and rejection review phases.
    \item The address of the assigned authority (\texttt{assignedAuthority}).
    \item The current lifecycle state (\texttt{ReportStatus} enum).
    \item The phase deadline timestamp (\texttt{phaseDeadline}).
    \item An append-only array of \texttt{AuthorityAction} entries, which records each authority interaction with a comment and an optional image CID.
\end{itemize}

To keep gas costs low, large data (images, long text) are not stored on-chain. Only the IPFS URI and cryptographic hashes are stored in the struct.

\subsubsection*{Role-Based Access Control}

The contract uses OpenZeppelin's \texttt{AccessControl} module to define two operational roles:

\begin{itemize}
    \item \texttt{RELAYER\_ROLE}: Assigned to the backend relayer wallet. Only this role is permitted to call \texttt{submitReport()}, \texttt{castValidationVote()}, \texttt{castVerificationVote()}, and \texttt{castRejectionReviewVote()}.
    \item \texttt{AUTHORITY\_ROLE}: Assigned to municipal authority workers through the multi-sig governance contract. Only this role is permitted to call \texttt{startWork()}, \texttt{addAuthorityUpdate()}, \texttt{markAsSolved()}, and \texttt{rejectIssue()}.
\end{itemize}

The ownership of \texttt{Reporting.sol} is transferred to \texttt{AuthorityMultiSig.sol} at deployment time, so only the governance contract can grant or revoke the \texttt{AUTHORITY\_ROLE}.

\subsubsection*{Eight-State Finite State Machine}

The report lifecycle is enforced by an eight-state \ac{FSM} as described in Chapter~\ref{ch:architecture}. The \texttt{\_finalizeSingleReport()} function is the central state transition function. It evaluates the current voting counts and the phase deadline to determine the correct next state.

\subsubsection*{Configurable Voting Method}

The voting outcome for all three community voting phases is determined by a configurable \texttt{VotingMethod} enum. The method is set by the contract owner using the \texttt{setVotingConfig()} function and applies globally to all reports. Four strategies are supported:

\begin{itemize}
    \item \textbf{Majority51} (default): A simple majority rule. The positive side wins if its vote count exceeds the negative count. No minimum participation is required.
    \item \textbf{SuperMajority66}: The positive side must receive at least two-thirds of the total cast votes. Formally: $\text{positive} \times 3 \geq \text{total} \times 2$.
    \item \textbf{Threshold}: A hard quorum must first be met ($\text{total} \geq \texttt{minVotesRequired}$). If the quorum is not met, the phase deadline is extended. Once the quorum is met, simple majority applies.
    \item \textbf{Hybrid}: Both of two configurable sub-methods (\texttt{h1} and \texttt{h2}) must independently pass for the positive side to win. Sub-methods cannot themselves be \texttt{Hybrid} to prevent infinite recursion.
\end{itemize}

The voting window duration is also configurable through \texttt{setVotingWindowDuration(duration)}. The default window is 6 hours. This allows the government authority to adjust the voting parameters to suit real-world consultation timelines without redeploying the contract.

\subsubsection*{Nullifier Tracking and Sybil Resistance}

Four persistent mappings track used ticket nullifiers to prevent Sybil attacks and replay attacks:

\begin{itemize}
    \item \texttt{usedSubmissionNullifiers}: Prevents a ticket from being used for more than one report submission.
    \item \texttt{usedValidationVoteNullifiers}: Prevents double voting in the community validation phase.
    \item \texttt{usedVerificationVoteNullifiers}: Prevents double voting in the post-resolution verification phase.
    \item \texttt{usedRejectionReviewVoteNullifiers}: Prevents double voting in the authority rejection appeal phase.
\end{itemize}

All state-changing functions follow the \ac{CEI} pattern: conditions are checked first, state mappings are updated next, and external interactions occur last. This pattern prevents reentrancy attacks.

\subsection{EmergencyReporting.sol}

\texttt{EmergencyReporting.sol} is a dedicated smart contract for civic emergency reports. Emergency reports represent situations that require immediate authority action without waiting for a community validation vote.

\subsubsection*{Emergency Report Data Structure}

Each emergency report is stored in an \texttt{EmergencyReport} struct containing the report ID, IPFS CID, report hash, submission nullifier, citizen pseudonym, relayer address, the current status, creation and update timestamps, the assigned authority address, the latest authority comment and image CID, and a reclassification flag.

\subsubsection*{Simplified Four-State Lifecycle}

Unlike the main \texttt{Reporting.sol} contract, \texttt{EmergencyReporting.sol} uses a simplified four-state lifecycle:
\begin{enumerate}
    \item \textbf{Open}: The report is submitted and immediately visible to authorities — no community validation vote is required.
    \item \textbf{InProgress}: An authority has claimed the report and started physical remediation.
    \item \textbf{Resolved}: The authority has completed the response.
    \item \textbf{Reclassified}: The authority has determined the report was not a genuine emergency (a false alarm).
\end{enumerate}

An authority can also resolve a report directly from the \texttt{Open} state, bypassing the \texttt{InProgress} step, to accommodate rapid responses.

\subsubsection*{False Alarm Penalty Box}

To prevent abuse of the emergency reporting channel, the contract implements a 30-day penalty box mechanism. If an authority calls \texttt{reclassifyEmergency(reportId, comment)}, the submitting citizen's pseudonym is recorded in the \texttt{emergencyPenaltyBox} mapping with an unlock timestamp set 30 days in the future:
\begin{equation}
    \texttt{penaltyUntil} = \texttt{block.timestamp} + 30 \times 24 \times 60 \times 60\,\text{s}
\end{equation}
Any subsequent emergency report submission by the same pseudonym before the penalty expires is rejected with a custom \texttt{EmergencyReportingLocked} error. After 30 days, the pseudonym is automatically unblocked and may submit emergency reports again.

\subsubsection*{Access Control}

\texttt{EmergencyReporting.sol} uses simple boolean mappings (\texttt{authorizedRelayers} and \texttt{authorizedAuthorities}) rather than the multi-sig governance of \texttt{Reporting.sol}. This allows the owner to quickly add or remove relayers and authority workers through the \texttt{setRelayer()} and \texttt{setAuthority()} functions.

\subsection{OpinionPolling.sol}

\texttt{OpinionPolling.sol} manages the public opinion polling workflow. It is used by authorized municipal officials to create polls on civic proposals and collect anonymous community responses.

Each poll is stored in a \texttt{Poll} struct containing the IPFS metadata CID, the deadline timestamp, the poll type (\texttt{TrueFalse} or \texttt{MultiChoice}), and a result mapping from option index to vote count. A separate nullifier mapping (\texttt{nullifierVoted[pollId][nullifier]}) is maintained to prevent double voting.

The contract verifies the caller's authority credentials against the \texttt{Reporting.sol} access registry before allowing poll creation. Votes are cast through the \texttt{castVote(pollId, optionIndex, nullifier)} function, which is called by the backend relayer on behalf of the citizen.

\subsection{AuthorityMultiSig.sol}

\texttt{AuthorityMultiSig.sol} governs all administrative actions within the system. It implements a two-tier multi-signature governance model:

\begin{itemize}
    \item \textbf{Super Admins}: Senior municipal officials (such as the Mayor or Chief Executive) who are authorized to submit and vote on governance proposals.
    \item \textbf{Authority Workers}: Municipal engineering staff who are authorized to act on civic reports.
\end{itemize}

Governance actions — including adding an authority (\texttt{AddAuthority}), removing an authority (\texttt{RemoveAuthority}), updating a Super Admin, and modifying the on-chain profile of an official — require a quorum of Super Admin votes. The quorum is computed dynamically as:
\begin{equation}
    \text{Quorum} = \Big\lfloor \frac{\texttt{superAdminCount}}{2} \Big\rfloor + 1
\end{equation}

Once a proposal reaches quorum, it is executed automatically. For \texttt{AddAuthority} and \texttt{RemoveAuthority} proposals, the contract makes a cross-contract call to \texttt{IReporting.setAuthority(target, authorized)}, which updates the \texttt{AUTHORITY\_ROLE} registry in \texttt{Reporting.sol}.

Each Super Admin and Authority Worker is associated with an immutable on-chain \texttt{Profile} struct containing their \texttt{name}, \texttt{position}, and \texttt{department}. Citizens can view these profiles through the DApp to verify which official is assigned to handle a reported civic issue.

\section{Identity Decoupling and ZKP GovID Simulator}
\label{sec:identity-impl}

The ZKP GovID Simulator is a Node.js/Express service that simulates a national government identity registry. It is responsible for authenticating citizens and issuing cryptographic credentials that are used for anonymous but verifiable interactions on the blockchain.

\subsection{Citizen Authentication}

Citizens authenticate by submitting their Government Identity Number (\texttt{govId}) and password to the \texttt{/api/govid/verify-citizen} endpoint. The service checks these credentials against a mock citizen database. This step simulates the interface between the system and a real national identity infrastructure.

\subsection{Citizen Seed and Keypair Derivation}

Upon successful authentication, the simulator generates a unique 256-bit \texttt{citizenSeed}. This seed is returned to the DApp. The DApp uses the \texttt{ethers.js} library to deterministically derive an anonymous ECDSA key pair from the seed. This key pair is stored only in the browser's local storage. The seed and the derived keys are never sent to any server, ensuring that the citizen's blockchain identity is known only to the citizen's own device.

\subsection{Ticket Batch Generation and Signing}

The simulator also generates a batch of ten one-time ticket identifiers (\texttt{zkpTicketId\textsubscript{1}} $\ldots$ \texttt{zkpTicketId\textsubscript{10}}). Each ticket is a cryptographic hash. The simulator signs each ticket using the government authority's ECDSA private key. These signed tickets are returned to the DApp along with the \texttt{citizenSeed}.

Each signed ticket functions as a government-endorsed authorization token. The citizen uses one ticket per report submission, one per vote, and so on. Once a ticket's nullifier has been recorded on-chain, it cannot be reused.

\section{Backend Relayer Middleware Implementation}
\label{sec:relayer-impl}

The backend relayer is implemented using the NestJS framework with TypeScript. It is the central middleware component of the system. All citizen-facing interactions with the blockchain and off-chain services are routed through the relayer.

\subsection{Dual-Layer Cryptographic Verification}

All incoming report payloads are verified by a dual-layer signature verification engine before any processing begins:

\begin{enumerate}
    \item \textbf{Government Ticket Verification}: The relayer fetches the government authority's ECDSA public key and verifies the ticket signature. If the signature is invalid, the request is rejected immediately. This confirms that the submitting citizen holds a valid government-issued credential.
    \item \textbf{Citizen Payload Integrity Verification}: The relayer reconstructs the composite report hash and verifies the citizen's ECDSA signature over it. This confirms that the report content has not been tampered with in transit.
\end{enumerate}

Only after both verifications pass does the relayer enqueue the report for processing.

\subsection{Asynchronous Processing with BullMQ Flow Queues}
\label{subsec:bullmq-impl}

After verification, the relayer does not process the report synchronously. Instead, it enqueues the report as an asynchronous BullMQ Flow Job in a Redis-backed queue and immediately returns a \texttt{jobId} to the DApp. The DApp uses this \texttt{jobId} to track progress via \ac{SSE}.

The queue module (\texttt{ReportQueueModule}) is implemented using the \texttt{@nestjs/bullmq} integration. Redis connection parameters (host, port, password) are configured through environment variables. Two queues are registered:

\begin{itemize}
    \item \texttt{report-processing}: The main queue for report flow parent jobs.
    \item \texttt{report-step-processing}: The step queue for the three child jobs.
\end{itemize}

The \texttt{ReportFlowProducer} service uses the BullMQ \texttt{FlowProducer} API to enqueue a nested job tree with the following structure:

% \begin{verbatim}
% Parent: Report [<category>] - rpt_<ticketHash>
%   └─ Step 3: step3-blockchain-submit (rpt_<ticketHash>-chain)
%        └─ Step 2: step2-ipfs-upload (rpt_<ticketHash>-ipfs)
%             └─ Step 1: step1-ai-moderation (rpt_<ticketHash>-ai)
% \end{verbatim}

BullMQ executes the tree bottom-up: Step 1 runs first, its result is read by Step 2, Step 2's result is read by Step 3, and finally the parent summary job runs. This dependency chain is enforced by BullMQ natively — a parent job is only moved to the active state after all its children have completed successfully.

Each child step returns a typed result object. These results are passed to the next step using the \texttt{job.getChildrenValues()} API:

\begin{itemize}
    \item Step 1 returns \texttt{AiStepResult}: \texttt{\{isApproved, blurredMedia\}}.
    \item Step 2 returns \texttt{IpfsStepResult}: \texttt{\{cid, ipfsUri\}}.
    \item Step 3 returns \texttt{BlockchainStepResult}: \texttt{\{transactionHash, blockNumber\}}.
\end{itemize}

Each step is configured with \texttt{failParentOnFailure: true} and up to three retry attempts using exponential backoff:
\begin{equation}
    \Delta t(k) = 2000 \times 4^{k-1} \text{ ms}, \quad k \in \{1, 2, 3\}
\end{equation}
This gives delays of 2\,s, 8\,s, and 32\,s between attempts. Step 1 (AI moderation) uses only one attempt, as an AI rejection is a deterministic content decision and not a transient failure.

An additional \textbf{Vote Queue} (\texttt{vote-queue}) handles community votes asynchronously. The \texttt{VoteQueueProcessor} receives vote job data (report ID, vote phase, nullifier, decision, pseudonym) and submits the vote transaction to the blockchain with up to three retry attempts.

\subsection{Real-Time Progress Notification (SSE Gateway)}
\label{subsec:sse-impl}

The \texttt{ReportStatusGateway} controller implements the \ac{SSE} notification layer. It uses an RxJS \texttt{Subject} as a push-based event bus. The processing flow is:

\begin{enumerate}
    \item At each pipeline step, the queue worker emits a \texttt{ReportJobProgress} object through the NestJS \texttt{EventEmitter2} bus.
    \item The \texttt{@OnEvent(REPORT\_JOB\_PROGRESS\_EVENT)} handler in \texttt{ReportStatusGateway} receives this event and calls \texttt{progressSubject.next(progress)}.
    \item Any DApp client connected to \texttt{GET /report/status/stream?citizenAddress=<address>} receives the event through an RxJS pipeline that filters by \texttt{citizenPubKey}.
    \item The DApp renders the received progress step and percentage in a live progress bar.
\end{enumerate}

The structured progress payload (\texttt{ReportJobProgress}) contains:

\begin{itemize}
    \item \texttt{jobId}: The root job identifier (e.g., \texttt{rpt\_a1b2c3d4}).
    \item \texttt{citizenPubKey}: The citizen's derived wallet address (used for routing).
    \item \texttt{step}: The current pipeline stage (e.g., \texttt{ai\_moderation\_passed}).
    \item \texttt{percent}: An integer from 0 to 100 indicating overall progress.
    \item \texttt{message}: A human-readable status description (e.g., ``AI moderation passed. Content approved.'').
    \item \texttt{data}: Optional structured payload, such as the IPFS CID or blockchain transaction hash.
\end{itemize}

\subsection{Authority Funding Service and Cron Orchestration}

The relayer also contains two automated background services:

\begin{enumerate}
    \item \textbf{AuthorityFundingService}: This service listens for \texttt{ProposalExecuted} events emitted by \texttt{AuthorityMultiSig.sol} via a WebSocket connection. When an \texttt{AddAuthority} proposal is executed, the service automatically transfers 5.0 ETH from the relayer treasury to the new Authority Worker's wallet. This ensures the worker has sufficient gas to submit blockchain transactions immediately after being granted access.
    \item \textbf{CronService}: Two scheduled jobs are run daily:
    \begin{itemize}
        \item \textit{Balance monitoring}: Every day at 01:00 (Colombo time zone), the service checks the ETH balances of all registered Super Admins and Authority Workers and replenishes any account below 1.0 ETH.
        \item \textit{Batch voting finalization}: Every day at midnight, the service calls \texttt{batchFinalizeVotingWindows()} on \texttt{Reporting.sol} for all reports in voting phases ($S_0$, $S_4$, $S_5$) whose \texttt{phaseDeadline} has elapsed. This prevents reports from remaining stuck in pending states.
    \end{itemize}
\end{enumerate}

\subsection{BullStudio Monitoring Dashboard}

A BullStudio monitoring dashboard is integrated into the relayer at the path \texttt{/ops/bullstudio}. It is protected by session-based authentication. The dashboard displays the status of jobs in both the \texttt{report-processing} and \texttt{report-step-processing} queues. It shows waiting, active, completed, and failed job counts, and provides a view of individual job details, logs, and retry history. This tool is used for system monitoring and debugging during testing and deployment.

\section{Off-Chain IPFS Storage Integration}
\label{sec:ipfs-impl}

The IPFS storage service is implemented as a dedicated NestJS module (\texttt{IpfsModule}). It provides an abstraction layer over the IPFS HTTP API.

\subsection{Service Architecture}

The IPFS service consists of two Docker containers:
\begin{enumerate}
    \item \textbf{IPFS Node}: A standard IPFS daemon (version 0.41.0) that manages the peer-to-peer content store.
    \item \textbf{IPFS HTTP Service}: A lightweight NestJS service (exposed on port 4000) that provides a REST API for file upload and retrieval. It communicates with the IPFS node through its RPC API on port 5001.
\end{enumerate}

\subsection{Upload Pipeline with AI Face Blurring}

The upload pipeline is integrated with the AI moderation step. After AI moderation is approved, the queue worker reads the \texttt{blurredMedia} array from the AI step result. If the AI oracle has detected and blurred any human faces in the submitted images, the blurred versions (returned as base64-encoded data) replace the original image buffers before the IPFS upload begins. This ensures that face data is not permanently stored in IPFS.

The \texttt{IpfsService.uploadComplaint()} method bundles the report description, category, location, and processed images into a complaint object and uploads it to IPFS. The returned \texttt{CID} and \texttt{ipfsUri} are stored in the step result and passed to the blockchain submission step.

\subsection{CID Integrity Verification}

The CID is cryptographically derived from the content of the uploaded file. The blockchain submission step stores this \texttt{ipfsUri} on-chain in the \texttt{Report} struct. At any time, a verifier can retrieve the content from IPFS by its CID and recompute the hash to confirm that the stored data has not been modified.

\section{AI Content Moderation Oracle Service Ensemble}
\label{sec:ai-impl}

The AI moderation service is implemented as a network of four containerized microservices. All four containers are deployed on a dedicated VPS and communicate through an internal Docker network.

\subsection{Oracle Worker Implementations and Model Hyperparameters}
\label{subsec:oracle-hyperparameters}

Each oracle worker uses pre-trained transformer or sentence-embedding models with tuned decision thresholds. The thresholds are set as environment variables so they can be adjusted without rebuilding the containers.

\subsubsection*{Safety Oracle (Oracle 1)}

The Safety Oracle uses two models:

\begin{enumerate}
    \item \textbf{Text Safety — \texttt{unitary/toxic-bert}}: A BERT-based multi-label classifier fine-tuned on the Jigsaw Toxic Comment dataset. It outputs six independent toxicity probability scores. Table~\ref{tab:safety-text-thresholds} shows the per-label decision thresholds used to classify text as unsafe.

    \begin{table}[h!]
    \centering
    \caption{Safety Oracle: unitary/toxic-bert Decision Thresholds}
    \label{tab:safety-text-thresholds}
    \begin{tabular}{|l|l|l|}
    \hline
    \textbf{Label} & \textbf{Threshold} & \textbf{Effect} \\
    \hline
    \texttt{toxic} & $\geq 0.50$ & Content flagged as unsafe \\
    \hline
    \texttt{severe\_toxic} & $\geq 0.50$ & Critical violation (override) \\
    \hline
    \texttt{obscene} & $\geq 0.60$ & Content flagged as unsafe \\
    \hline
    \texttt{threat} & $\geq 0.40$ & Critical violation (override) \\
    \hline
    \texttt{insult} & $\geq 0.60$ & Content flagged as unsafe \\
    \hline
    \texttt{identity\_hate} & $\geq 0.40$ & Critical violation (override) \\
    \hline
    \end{tabular}
    \end{table}

    Text is truncated to 512 tokens before inference to match the model's maximum input length. A \textit{critical violation} from \texttt{severe\_toxic}, \texttt{threat}, or \texttt{identity\_hate} triggers a direct aggregator-level \texttt{REJECT} override.

    \item \textbf{Image Safety — \texttt{Falconsai/nsfw\_image\_detection}}: An image classification model that outputs probability scores for the labels \texttt{nsfw}, \texttt{porn}, \texttt{sexy}, and \texttt{hentai}. The maximum score across these labels is taken as the composite NSFW score. An image is classified as unsafe if its NSFW score is $\geq 0.70$.

    \item \textbf{Face Blurring — OpenCV Haar Cascades}: Safe images are processed for face detection using four OpenCV Haar cascade classifiers: \texttt{haarcascade\_frontalface\_default.xml}, \texttt{haarcascade\_frontalface\_alt.xml}, \texttt{haarcascade\_frontalface\_alt2.xml}, and \texttt{haarcascade\_profileface.xml} (applied twice: on the original and horizontally flipped image for right-profile detection). Face detection parameters are: \texttt{scaleFactor = 1.05}, \texttt{minNeighbors = 3}, \texttt{minSize = (20, 20)} pixels. Overlapping detections are deduplicated using Non-Maximum Suppression (NMS) with an overlap threshold of 0.4. Detected face regions are blurred with a Gaussian kernel sized proportionally to the face bounding box, with a minimum kernel size of $19 \times 19$.
\end{enumerate}

\subsubsection*{Spam and Abuse Oracle (Oracle 2)}

The Spam and Abuse Oracle combines an NLP model with a rule-based scoring engine:

\begin{enumerate}
    \item \textbf{NLP Model — \texttt{mrm8488/bert-tiny-finetuned-sms-spam-detection}}: A lightweight BERT-tiny model fine-tuned on SMS spam data. It is used as a signal alongside rule-based checks.

    \item \textbf{Decision Thresholds}: Table~\ref{tab:spam-thresholds} shows the four configurable thresholds used to combine the AI model score and the rule-based score into a final spam decision.

    \begin{table}[h!]
    \centering
    \caption{Spam Oracle: Hybrid Decision Thresholds}
    \label{tab:spam-thresholds}
    \begin{tabular}{|l|l|l|}
    \hline
    \textbf{Parameter} & \textbf{Value} & \textbf{Meaning} \\
    \hline
    \texttt{AI\_STRONG\_SPAM\_THRESHOLD} & 0.85 & AI score alone is sufficient to reject \\
    \hline
    \texttt{AI\_SUPPORTING\_SPAM\_THRESHOLD} & 0.65 & AI score supports rejection if rules agree \\
    \hline
    \texttt{RULE\_REJECT\_THRESHOLD} & 0.60 & Rule score alone is sufficient to reject \\
    \hline
    \texttt{RULE\_SUPPORT\_THRESHOLD} & 0.25 & Rule score supports rejection if AI agrees \\
    \hline
    \end{tabular}
    \end{table}

    \item \textbf{Rule-based signals} and their contribution to the spam score are: a text length below 5 words adds 0.35; a matched spam phrase adds 0.80; promotional language adds 0.35; more than one URL adds 0.60; one URL adds 0.25; repeated character patterns add 0.50; excessive symbols add 0.40. Scores are capped at 1.0.
\end{enumerate}

\subsubsection*{Civic Relevance Oracle (Oracle 3)}

\begin{enumerate}
    \item \textbf{Embedding Model — \texttt{sentence-transformers/all-MiniLM-L6-v2}}: A compact sentence transformer model that maps text to 384-dimensional semantic embedding vectors. Civic reference descriptions for 9 issue categories are pre-embedded at startup and stored in memory.

    \item \textbf{Decision Thresholds}: Table~\ref{tab:civic-thresholds} shows the two configurable similarity thresholds.

    \begin{table}[h!]
    \centering
    \caption{Civic Relevance Oracle: Cosine Similarity Thresholds}
    \label{tab:civic-thresholds}
    \begin{tabular}{|l|l|l|}
    \hline
    \textbf{Parameter} & \textbf{Value} & \textbf{Meaning} \\
    \hline
    \texttt{TEXT\_RELEVANCE\_THRESHOLD} & 0.38 & Minimum cosine similarity to a reference civic description \\
    \hline
    \texttt{CATEGORY\_MATCH\_THRESHOLD} & 0.40 & Minimum similarity for declared category to match text \\
    \hline
    \end{tabular}
    \end{table}

    The cosine similarity is computed between the report embedding and each reference embedding for the nine civic categories: Road Damage, Waste Management, Streetlight Issue, Drainage/Sewage, Water Supply, Flooding, Public Property Damage, Traffic/Road Safety, and Environmental Issues. A report is accepted if its maximum cosine similarity across all reference descriptions exceeds \texttt{TEXT\_RELEVANCE\_THRESHOLD = 0.38}. The category match threshold applies separately and does not cause rejection; it only adds a warning to the response.
\end{enumerate}

\subsection{Oracle Aggregator}
The Oracle Aggregator is the single entry point for the backend relayer. It:
\begin{enumerate}
    \item Receives the signed moderation request from the relayer. The relayer signs the request with its private key, and the aggregator recovers and verifies the signer address.
    \item Validates the request timestamp (to prevent stale requests) and nonce (to prevent replay attacks).
    \item Forwards the moderation request to all three oracle workers in parallel.
    \item Applies the following decision logic:
    \begin{itemize}
        \item If the Safety Oracle detects a critical violation, the final decision is \texttt{REJECT} regardless of other votes.
        \item If the Civic Relevance Oracle determines the content is non-civic, the final decision is \texttt{REJECT}.
        \item Otherwise, majority voting is applied: if two or more oracle workers return \texttt{ACCEPT}, the final decision is \texttt{ACCEPT}; otherwise it is \texttt{REJECT}.
    \end{itemize}
    \item Computes a \texttt{report\_hash}, a \texttt{decision\_hash}, and an \texttt{aggregator\_signature} over the final decision. These fields are returned to the relayer in the moderation response.
\end{enumerate}

\section{Web-Based Decentralized Application (DApp)}

\subsection{Authentication Interface}
The login page collects the citizen's GovID credentials and sends them to the ZKP GovID Simulator. Upon successful verification, the DApp receives the \texttt{citizenSeed} and the signed ticket batch. The \texttt{ethers.js} library derives the anonymous ECDSA key pair from the seed. A custom React Context (\texttt{CitizenContext}) is used to store the keys and ticket batch in the browser's local storage. Keys and tickets never leave the device.

\subsection{Report Submission Interface}
The report submission form uses a mobile-first layout. Citizens enter a description, select a category, and optionally attach photographic evidence. On submission:
\begin{enumerate}
    \item The DApp hashes the description and ticket ID using Keccak-256.
    \item The citizen's private key signs the hash to produce the payload signature.
    \item The signed payload, images, and ZKP ticket are sent to the backend relayer via a multipart HTTPS POST request.
    \item The relayer responds immediately with a \texttt{jobId}.
    \item The DApp opens an \ac{SSE} connection to \texttt{/report/status/stream?citizenAddress=<address>} and displays a live progress bar.
\end{enumerate}

\subsection{Real-Time Submission Progress UI}
The progress UI is driven by SSE events received from the backend relayer. The DApp renders a step-by-step progress bar with the following stages:
\begin{enumerate}
    \item \textit{Queued}: Report is accepted and waiting for processing.
    \item \textit{AI Moderation}: Content is being analyzed by the oracle ensemble.
    \item \textit{IPFS Upload}: Approved content is being stored off-chain.
    \item \textit{On-Chain Submission}: Transaction is being submitted to the blockchain.
    \item \textit{Completed}: Transaction is confirmed on-chain. The transaction hash and block number are displayed.
\end{enumerate}
If any step fails, the UI displays a clear error message with the reason returned by the relayer (e.g., ``Content was rejected by AI moderation: non-civic content detected.'').

\subsection{Community Feed Interface}
The public feed retrieves submitted reports from the blockchain and displays them as interactive cards. Each card shows the report category, description excerpt, current status badge (e.g., \texttt{PENDING VALIDATION}, \texttt{OPEN}, \texttt{IN PROGRESS}, \texttt{CLOSED}), upvote and downvote counts, timestamp, and a thumbnail of the attached image retrieved from IPFS.

Citizens can filter the feed by category. Authenticated citizens can cast validation votes, verification votes, and rejection review votes directly from the feed. Vote submissions are also processed asynchronously through the vote queue, and the DApp receives SSE progress events for votes in the same way as for reports.

\subsection{Authority Dashboard}
Authenticated authority workers access a separate dashboard that shows all reports in the \texttt{Open} and \texttt{Reopened} states. Authorities can claim a report by calling \texttt{startWork()}, post mid-repair progress updates using \texttt{addAuthorityUpdate()}, mark a report as solved using \texttt{markAsSolved()}, or reject an invalid report using \texttt{rejectIssue()}. Each action triggers a blockchain transaction signed by the authority's wallet.

Super Admin users access the governance dashboard to submit and vote on multi-sig proposals for adding or removing authority workers.

\section{Deployment Infrastructure}
\label{sec:deployment-impl}

The complete system is deployed across two VPS instances to separate the core operational services from the resource-intensive AI and storage services.

\subsection{Server Allocation}

\textbf{Server 1 (Core Operations)} hosts:
\begin{itemize}
    \item Private PoA blockchain (3 validator nodes + 1 RPC node, via Docker Compose).
    \item ZKP GovID Simulator (Node.js/Express, Docker container).
    \item Backend Relayer with Redis (NestJS, Docker container).
    \item Web DApp (Next.js, Docker container).
\end{itemize}

\textbf{Server 2 (AI Oracle and Storage)} hosts:
\begin{itemize}
    \item AI Oracle Aggregator (Python/FastAPI, Docker container).
    \item Safety Oracle worker (Docker container).
    \item Spam and Abuse Oracle worker (Docker container).
    \item Civic Relevance Oracle worker (Docker container).
    \item IPFS Node and IPFS HTTP Service (Docker containers).
\end{itemize}

\subsection{CI/CD Pipeline}

Continuous integration and deployment are managed using Dokploy, a self-hosted Platform-as-a-Service tool installed on both servers. GitHub webhook triggers are configured on the \texttt{develop} branch of the project repository. When a commit is pushed to the branch, Dokploy automatically pulls the latest code, rebuilds the affected Docker images, and redeploys the updated containers. This ensures consistent and reproducible deployments across all environments.

\begin{figure}[!htbp]
    \centering
    \fbox{\parbox{0.85\textwidth}{\centering\vspace{1.5cm}
    \textit{[Image suggestion IMG-06: 2-Server VPS Deployment Topology Diagram. Show Server 1 hosting blockchain nodes, Redis, relayer, DApp, ZKP simulator. Show Server 2 hosting 4 AI oracle containers and 2 IPFS containers. Indicate Dokploy CI/CD on both servers and GitHub webhook trigger.]}
    \vspace{1.5cm}}}
    \caption{Two-Server VPS Deployment Topology}
    \label{fig:deployment-topology}
\end{figure}
